% Subsection 4.3.2: 超参数与训练策略

三种适配器共用同一套超参数，仅在样本构造层面不同（4.3.1 节）。LoRA 结构上取秩 $r=32$、缩放系数 $\alpha=64$（有效学习率 $\alpha/r=2$），并将适配层注入 ChatGLM 块中的 attention（融合 QKV 与输出）与 MLP 两侧线性层共四类位置，使 CoT 推理段与代码段在参数侧均可获得适配空间——这一点在 \texttt{LoRA\_cot\_*} 的监督目标下尤为关键。优化器侧采用 AdamW 配 cosine 调度（含约 3\% 步的线性 warmup）、bf16 数值精度与 gradient\_checkpointing 以匹配 9B 基座 + 32GB 单卡的显存预算；训练 3 个 epoch、等效 batch size 为 32，与训练集规模（约 9000 条）匹配，足以让 LoRA 收敛而不在 CoT 监督目标上过拟合。具体超参取值在仓库 \texttt{SETUP.md} 中以表格形式给出。

\subsubsection{设计依据与一致性约束}
上述取值遵循以下三条原则，使得 10 组实验之间仅因训练样本构造差异而产生可解释的指标变化：
\begin{itemize}
    \item \textbf{跨适配器一致}：三种 LoRA 共享完全相同的结构与优化器超参，仅在样本构造与示例条数 $k$ 上有差异；
    \item \textbf{与推理侧一致}：\texttt{LoRA\_code} 路径的 chat 模板与 4.2.1 节 Direct 推理对齐；\texttt{LoRA\_cot\_*} 路径的用户指令与 assistant 形态与 4.4 节 CoT 推理（含 few-shot 回放）对齐，避免特殊 token 或围栏格式在 train/test 间错位；
    \item \textbf{与数据规模匹配}：3 epoch 与等效 batch size 在 KodCode 训练集上约合数百次 optimizer step，足以让 LoRA 收敛而不至于过拟合。
\end{itemize}
